% Section 2: what musicians earn, from the microdata. This is the report's
% quantitative spine, because it is the only current source that sees the
% self-employed majority of this workforce.
\section{What Austin musicians earn}
\label{sec:earnings}

The Census Bureau's American Community Survey (ACS) is the one federal
instrument that interviews households directly, records the occupation of the
person answering, and asks about self-employment income alongside wages and
salaries.\footnote{The specific extract used throughout this section is the
Census Bureau's ACS 2020--2024 five-year Public Use Microdata Sample (PUMS) for
Texas, released 2026-03-05 and downloaded from
\url{https://www2.census.gov/programs-surveys/acs/data/pums/2024/5-Year/csv_ptx.zip}.
The PUMS is the individual-record file that lets a researcher recompute any
tabulation the Bureau publishes and many it does not; a parallel housing file
supplies the household records used for the rent statistics later in this
section. The Bureau samples roughly one percent of U.S.\ households each year
and pools five successive samples to make small-cell estimates like the
Austin-musician cell publishable, so the effective sample is the 2020, 2021,
2022, 2023 and 2024 waves stacked into one file. Musicians and singers are
Census occupation code (OCCP) 2752; music directors and composers are OCCP
2751. The Austin metro throughout this section is the five-county
Austin-Round Rock-San Marcos Metropolitan Statistical Area (Bastrop, Caldwell,
Hays, Travis and Williamson counties), mapped from the survey's PUMA
identifiers with the 2020 Census tract-to-PUMA relationship file. Employed
status uses the Bureau's employment recode ESR in \{1, 2, 4, 5\} (civilian at
work, civilian with a job but not at work, armed forces at work, or armed
forces with a job but not at work).} That combination is what payroll surveys
and administrative wage records cannot deliver: a payroll instrument sees only
workers on someone else's payroll, and administrative wage records see only
the covered universe of a particular tax or benefits program. Since
Section~\ref{sec:payroll} shows that most Austin musicians work for
themselves, the ACS is the one instrument whose universe includes them at all.
Every estimate below pools five years of Texas
microdata\footnote{We put dollar amounts on one basis in two steps, and the
order matters. For person \(i\) the real value is
\(y_i = Y_i \times (\mathrm{ADJINC}_i / 10^{6}) \times d_{2024}\), where \(Y_i\)
is the amount as reported, \(\mathrm{ADJINC}_i\) is the survey's own income
adjustment factor for that person's survey year, which restates all five
survey years on a common 2024 basis and is published with six implied
decimals, which is what dividing by \(10^{6}\) removes. The second factor is
\(d_{2024} = C_{2025}/C_{2024}\), the ratio of Bureau of Labor Statistics
Consumer Price Index for All Urban Consumers (CPI-U) annual averages, equal to
\NPumsDeflTwoZeroTwoFourToTwoZeroTwoFive, and it carries 2024 dollars to this
study's 2025 base. Housing dollars use the survey's housing factor, ADJHSG,
in the same position. There is no merge on survey year and no year-specific
deflator, because the survey factor has already removed the inflation inside
the 2020 to 2024 window and applying a second one would remove it twice.} and
uses one of two population definitions, each carried through to both sides of
every comparison it enters. We compute the medians and the earnings regression on
the positive-earnings base, employed adults aged 18 to 64 with personal
earnings above zero. We compute the self-employment share, the health-coverage
rates, the hours and weeks figures, and the earnings-threshold shares on the
employed base, which drops the positive-earnings requirement so that people
reporting zero or negative earnings from work stay in the denominator. Each
figure below names the base it uses.\footnote{Matching the base on both sides
matters more here than in most analyses. Comparing musicians drawn from an
unrestricted file, one that includes people out of the labour force and
people reporting zero earnings, against a baseline restricted to employed
workers moves the musician median from about \$21,000 to about \$12,800 on
the mismatch alone. Every estimate in this section uses matched bases and
reports the unweighted record count behind it.} Margins of error come from
the Census Bureau's replicate-weight method.\footnote{The Bureau's published
variance method for this file is successive-difference replication, and it
is the method it recommends for any statistic computed from the PUMS. We
recompute every statistic on each of the 80 replicate weightings the Bureau
distributes, and the spread of those 80 answers around the full-sample answer
is the sampling variance:
\[
\widehat{\operatorname{Var}}(\hat\theta)
  = \frac{4}{80}\sum_{r=1}^{80}\bigl(\hat\theta_r - \hat\theta\bigr)^{2},
\]
where \(\hat\theta\) is the estimate computed with the full person weight
PWGTP, \(\hat\theta_r\) is the same estimate computed with replicate weight
PWGTP\(r\), and the constant \(4/80\) is fixed by the Bureau's replicate
design rather than chosen here. We take deviations around the full-sample
estimate rather than around the average of the 80 replicates, the convention
Stata labels \texttt{mse}, so the uncertainty on medians, means and
regression coefficients in this section is directly comparable across
statistics. The reported margin is
\(\mathrm{MOE}_{90} = 1.645\,\sqrt{\widehat{\operatorname{Var}}(\hat\theta)}\),
where 1.645 is the standard normal value that leaves 5 percent in each tail.
Treating the records as an independent sample would understate the error
substantially, because the Bureau's sample is clustered and its weights are
calibrated to independent population controls. A small number of replicate
weight cells are slightly negative, which the Bureau permits and which is an
artefact of how the replicates are built; a weighted percentile is undefined
with negative weights, so we floor the weight at zero inside the percentile
calculation only, and means, shares and regressions use the replicate weights
as published.}

\subsection{Austin musicians earn about half what other Austin workers earn}

Austin musicians report median personal earnings from work of
\NPumsAustinMusMedpernpPos\ (\NPumsAustinMusMedpernpPosMoe) against
\NPumsAustinAllworkMedpernpPos\ (\NPumsAustinAllworkMedpernpPosMoe) for all
Austin workers on the same definition.\footnote{The universe on both sides is
employed adults aged 18 to 64 living in the Austin-Round Rock-San Marcos MSA
with positive personal earnings from work, drawn from the ACS 2020--2024
five-year Texas PUMS. The musician cell uses OCCP 2752 (Musicians and
Singers) and comes from \NPumsAustinMusNPos\ unweighted person records
representing a weighted population of \NPumsAustinMusWtdPos\ people; the
all-workers cell comes from 48,606 unweighted records representing 1.28 million
people. Both figures use the Bureau's person weight PWGTP, so each record
represents the number of people its weight says it represents. A weighted
median is the smallest observed value at which the weight accumulated at or
below that value reaches half of the total weight,
\(\hat\theta_{0.5} = \inf\{\,y : \sum_i w_i\,\mathbf{1}(y_i \le y)
   \ge 0.5\sum_i w_i\,\}\),
where \(y_i\) is record \(i\)'s earnings in 2025 dollars, \(w_i\) is PWGTP,
and \(\mathbf{1}(\cdot)\) is one when the condition inside holds and zero
otherwise; adjacent order statistics are averaged when the accumulated weight
lands exactly on half. A weighted share of some group \(G\) is
\(\hat p = \sum_{i \in G} w_i a_i \big/ \sum_{i \in G} w_i\), with \(a_i\)
equal to one when the record has the attribute being counted, so the
denominator is the weighted size of the group and not the record count. Both
are descriptive population estimates: they describe who the survey
enumerated, and no statistical model stands behind them.} The estimation
sample for the Austin musician median is \NPumsAustinMusNPos\ unweighted
records over five survey years, enough to publish and not enough to break
down further, so this study reports no Austin musician estimates by race,
age or sex.

\begin{figure}[htbp]
\centering
\includegraphics[width=\textwidth]{fig02_earnings_gap.png}
\caption{Median personal earnings from work and median total personal income by
occupation, Texas workers aged 18 to 64 with positive earnings, 2025 dollars.
Source: American Community Survey 2020--2024 five-year public use microdata,
person-weighted; whiskers are 90 percent margins of error from the 80 replicate
weights. How to read it: the orange bar is what the work pays and the navy bar
is total personal income from all sources, so the gap between them is income
from somewhere other than the person's own work. Caveats:
the survey records one occupation per person, so a musician who reports a day
job as their main work is counted in that occupation instead; the combined
media, audio-visual and sound-equipment category is a single Census code, not
sound engineers alone.}
\label{fig:earningsgap}
\end{figure}

Figure~\ref{fig:earningsgap} covers Texas, and what matters in it is less the
ranking than the distance between each occupation's two bars. Actors show the
widest pair, with total personal income roughly twice their earnings from
work; musicians show a smaller but visible pair, and all Texas workers show
almost none. The Austin cell shows the same pattern. Austin musicians report
median total personal income of \NPumsAustinMusMedpincpPos\ against
\NPumsAustinMusMedpernpPos\ in earnings from work. For all Austin workers the
same two figures are \NPumsAustinAllworkMedpincpPos\ and
\NPumsAustinAllworkMedpernpPos, a difference of a few hundred dollars.
The distance between two separately estimated medians is not any one person's
non-work income, but it does say that Austin musicians as a group draw on
income beyond their own work. The survey records income by source in seven
separate fields (wages and salary, self-employment, interest and dividends,
retirement, Social Security, Supplemental Security Income, and public
assistance), and decomposing the gap into those fields would be a direct
extension of the work here.

Hours worked do not account for the low earnings either. The median Austin
musician reports a usual week of \NPumsAustinMusMedwkhp\ hours (the survey
variable WKHP) and \NPumsAustinMusMedwkwn\ weeks worked in the past year
(WKWN), matching \NPumsAustinAllworkMedwkhp\ hours and
\NPumsAustinAllworkMedwkwn\ weeks for the median Austin worker. Both medians
come from the same 77-record Austin cell and both equal the modal answers,
so they say little about the tails: statewide, where the sample is larger,
the median musician's usual week is \NPumsTxMusMedwkhp\ hours against
\NPumsTxAllworkMedwkhp\ for all Texas workers. Firmer evidence comes from
the regression below, which holds hours and weeks fixed and still leaves a
difference.

To probe whether the raw gap comes from age, education, hours or geography
rather than from something specific to music, we estimate a weighted least
squares regression of log personal earnings on an indicator for music
occupations, with controls for age, sex, education, usual hours and weeks
worked, and an indicator for each metro area.\footnote{The fitted equation is
\[
\ln y_i = \alpha + \beta\,\mathrm{music}_i
  + \gamma_1 \mathrm{age}_i + \gamma_2 \mathrm{age}_i^{2}
  + \gamma_3 \mathrm{female}_i
  + \sum_{k=2}^{5}\delta_k\,\mathrm{educ}_{ki}
  + \gamma_4 \mathrm{hours}_i + \gamma_5 \mathrm{weeks}_i
  + \sum_{m}\lambda_m\,\mathrm{metro}_{mi} + \varepsilon_i .
\]
The dependent variable \(y_i\) is person \(i\)'s personal earnings from work
in 2025 dollars (the ACS PERNP field, deflated as described above).
\(\mathrm{music}_i\) is one when the person's occupation is musician, singer,
music director or composer (OCCP 2751 or 2752) and zero otherwise;
\(\mathrm{age}_i\) is AGEP, age in years; \(\mathrm{female}_i\) is one for
women (SEX = 2); the \(\mathrm{educ}_{ki}\) are indicators for four
educational-attainment bands drawn from SCHL, measured against a base of
less than high school (band 2 is high-school diploma or GED, band 3 is some
college, band 4 is bachelor's degree, band 5 is graduate or professional
degree); \(\mathrm{hours}_i\) is WKHP (usual hours worked per week);
\(\mathrm{weeks}_i\) is WKWN (weeks worked in the past twelve months); and
the \(\mathrm{metro}_{mi}\) are Metropolitan Statistical Area indicators
with the Austin-Round Rock-San Marcos MSA as the omitted reference, so each
reported metro coefficient is a difference from Austin and the balance of
the state enters as its own estimated term rather than as the base.
\(\varepsilon_i\) is the residual. Estimation is by weighted least squares
with the person weight PWGTP over the \NPumsRegN\ employed Texans aged 18 to
64 with positive earnings, and the standard error comes from the same 80
replicate weights used for the medians, the Census Bureau's own
successive-difference method, rather than from an assumption that records
are independent. Because the dependent variable is in logs, a coefficient
converts to a percentage difference as \(100\,(e^{\hat\beta} - 1)\), and the
same transformation applied to \(\hat\beta \pm 1.96\,\mathrm{se}(\hat\beta)\)
gives the 95 percent interval. The music coefficient
\NPumsRegMusicCoef\ with standard error \NPumsRegMusicSe\ puts the 95 percent
interval at \(-32.7\) to \(-16.8\) percent. Two things change at once between
this estimate and the Austin medians above: the population is every employed
Texan aged 18 to 64 with positive earnings rather than the Austin metro cell,
and the occupation is musicians, singers, music directors and composers (OCCP
2751 and 2752) rather than musicians and singers alone (OCCP 2752). This is a
descriptive conditional difference, not a causal estimate of what choosing
music does to earnings: people who choose music differ from people who do
not in ways the survey does not record, including talent, family resources
and preference for the work itself. The full three-column table (metro fixed
effects only, plus demographics, plus hours and weeks) is at
\texttt{03\_analysis/out/tables/table\_earnings\_regression.csv}, and the
figures reported in this paragraph come from column (3).}
Among employed Texans aged 18 to 64 with positive earnings from work,
adjusting for age, sex, education, usual hours, weeks worked and metro area,
music occupations show a \NPumsRegMusicPctgap\ earnings difference, against a
raw Austin shortfall of about half. Part of the raw gap therefore tracks
those characteristics and part does not.

\subsection{Most of this workforce works for itself}

\begin{figure}[htbp]
\centering
\includegraphics[width=0.86\textwidth]{fig03_selfemp_share.png}
\caption{Share of workers who are self-employed, incorporated or not, by
occupation. Source: American Community Survey 2020--2024 five-year microdata,
employed Texans aged 18 to 64; whiskers are 90 percent margins of error. How to
read it: each bar is the share of workers in that occupation whose main job is
their own business rather than someone else's payroll. Caveat: the survey
records class of worker for the main job only, so a musician who also holds a
payroll day job is counted under the day job.}
\label{fig:selfemp}
\end{figure}

\NPumsAustinMusSelfemp\ (\NPumsAustinMusSelfempMoe) of Austin musicians work
for themselves, against \NPumsAustinAllworkSelfemp\
(\NPumsAustinAllworkSelfempMoe) of all Austin workers, which makes an Austin
musician about seven times as likely to be self-employed as the typical Austin
worker.\footnote{Universe: employed adults aged 18 to 64 in the
Austin-Round Rock-San Marcos MSA, ACS 2020--2024 five-year Texas PUMS. The
numerator on both sides is class-of-worker categories 6 (self-employed,
unincorporated) and 7 (self-employed, incorporated) from the survey's COW
field; the denominator is the employed base on that geography. The musician
share comes from \NPumsAustinMusNEmp\ unweighted person records, so we do not
break it down by demographic subgroup.} Statewide the musician rate is lower,
\NPumsTxMusSelfemp\ (\NPumsTxMusSelfempMoe) against \NPumsTxAllworkSelfemp\
for all Texas workers, so the Austin estimate is well above the state
estimate on a cell of only \NPumsAustinMusNEmp\ records and should be read
with that in mind.

Self-employment at that rate accounts for much of the distance between the
payroll count and the survey count, though not all of it. Applying the Austin
self-employment share to the \NPumsAustinMusWtdPos\ musicians the survey
estimates leaves roughly 600 on someone's payroll, against the
\NEmpPayrollMusiciansAustinLast\ jobs the federal wage survey recorded in the
Austin metro in 2025.\footnote{The federal wage survey here is the Bureau of
Labor Statistics Occupational Employment and Wage Statistics (OEWS) program,
a semi-annual establishment survey administered by BLS and the state
workforce agencies. Program documentation:
\url{https://www.bls.gov/oes/oes_ques.htm}. OEWS reports the wage of every
worker on a sampled establishment's payroll and reports nothing for the
self-employed, who fall outside its sampling frame; that is the arithmetic
reason a survey which counts self-employment income and a payroll survey
which does not disagree so widely on the same occupation and the same metro.
Section~\ref{sec:payroll} takes up the payroll count separately.} It also
explains why a payroll-based wage series can look healthy while a survey
which counts self-employment income does not.

\subsection{The distribution matters more than the median}

\begin{figure}[htbp]
\centering
\includegraphics[width=0.86\textwidth]{fig04_earnings_distribution.png}
\caption{Cumulative share of workers earning below each level, musicians against
all workers, Texas, 2025 dollars. Source: American Community Survey 2020--2024
five-year microdata, employed adults 18 to 64 with positive earnings. How to
read it: the height of a line at a given earnings level is the share of that
group earning less than that amount, so a line sitting higher and further left
means more of the group earns little. Caveat: earnings here come from all work,
not from music alone, so a musician with a day job appears at the combined
figure.}
\label{fig:earningsdist}
\end{figure}

The spread below the median is wide. Figure~\ref{fig:earningsdist} plots
Texas, where the sample is large enough to draw a smooth curve; the shares
below are for the Austin cell, which comes from 77 records and so has
margins of error several times wider.
\NPumsAustinMusLtOneZerok\ (\NPumsAustinMusLtOneZerokMoe) of employed Austin
musicians aged 18 to 64 earn under \$10,000 from work, against
\NPumsAustinAllworkLtOneZerok\ (\NPumsAustinAllworkLtOneZerokMoe) of all
employed Austin workers on the same base;
\NPumsAustinMusLtTwoFivek\ (\NPumsAustinMusLtTwoFivekMoe) earn under
\$25,000 against \NPumsAustinAllworkLtTwoFivek; and
\NPumsAustinMusLtFiveZerok\ (\NPumsAustinMusLtFiveZerokMoe) earn under
\$50,000 against \NPumsAustinAllworkLtFiveZerok. Measured against the
\$10,000 line the 2014 Austin Music Census drew for 2013 music income,
restated to today's prices,\footnote{The 2013 threshold is the \$10,000 line the
2014 Austin Music Census applied to its respondents' 2013 pre-tax music
income. Titan Music Group, LLC, for the City of Austin Economic Development
Department, Music and Entertainment Division, \textit{The Austin Music
Census: A Data-Driven Assessment of Austin's Commercial Music Economy}
(fielded late 2014; published June 1, 2015; income data are pre-tax calendar
2013), \url{https://austin.widen.net/content/om7zqn1qbx/pdf/Austin_Music_Census_Interactive_PDF_53115.pdf}.
That census was a self-selected online convenience sample of 3,968
respondents drawn from the Austin music community; its published shares
describe the people who answered, and no probability-sample margin of error
attaches to them. Its microdata were never publicly released.}
\NPumsAustinMusLtmusiccensus\ (\NPumsAustinMusLtmusiccensusMoe) of Austin
musicians fall below the line that census drew, against
\NPumsAustinAllworkLtmusiccensus\ of the wider Austin workforce. That share
is not comparable with the 2014 census's own 68.4 percent: the 2014 census
measured music income only, among self-selected respondents, while the ACS
figure measures earnings from all work on a probability sample eight years
later.\footnote{Each share is the weighted proportion defined above,
\(\hat p = \sum_i w_i\,\mathbf{1}(y_i < \bar y)\big/\sum_i w_i\), where
\(\bar y\) is the threshold and \(y_i\) is personal earnings in 2025 dollars.
These four shares use the employed base (employed Austin musicians aged 18
to 64), rather than the positive-earnings base used for the medians above:
people reporting zero or negative earnings from work stay in the
denominator, because dropping them would remove the worst-off cases from a
statistic whose subject is how many people fall below a low threshold. The
2013 threshold is carried forward as
\(\$10{,}000 \times d_{2013}\), where \(d_{2013} = C_{2025}/C_{2013}\) is the
ratio of CPI-U annual averages, equal to
\NPumsDeflTwoZeroOneThreeToTwoZeroTwoFive, which puts the line at
\NPumsMusiccensusTenkTwoZeroTwoFiveusd\ in today's prices.}

\subsection{Nothing suggests Austin pays musicians better than other Texas metros}

Austin's median-musician earnings are lower than the medians in Dallas-Fort
Worth and Houston, but only by amounts small enough that the sampling error
in each metro cell easily covers the difference. These three metro figures
pool musicians, singers, music directors and composers (OCCP 2751 and 2752),
a wider grouping than the musicians-and-singers headline above, which is
what makes the metro cells large enough to publish.\footnote{The pooled
sample sizes are: Austin \NPumsMetroAustinMusicN\ unweighted records
representing \NPumsMetroAustinMusicWtd\ people; Dallas-Fort Worth 191 records
representing 4,505 people; Houston 117 records representing 3,453 people.
Each cell is drawn from the ACS 2020--2024 five-year Texas PUMS, employed
adults aged 18 to 64 with positive earnings, and each metro is its full
five- or ten-county MSA as defined by OMB's 2023 core-based statistical area
delineation. Musicians and singers alone (OCCP 2752) yield 33 Dallas records
and 26 Houston records, both under the 50-record suppression threshold this
study uses.} Median personal earnings from work in these pooled music
occupations are \NPumsMetroAustinMusicMedpernp\
(\NPumsMetroAustinMusicMedpernpMoe) in Austin,
\NPumsMetroDfwMusicMedpernp\ (\NPumsMetroDfwMusicMedpernpMoe) in Dallas-Fort
Worth and \NPumsMetroHoustonMusicMedpernp\
(\NPumsMetroHoustonMusicMedpernpMoe) in Houston. Austin has the lowest point
estimate of the three, and every margin is wide, so neither pairwise gap is
statistically firm.\footnote{Austin against Houston is a difference of about
\$12,500 against a combined 90 percent margin of roughly \$18,600; Austin
against Dallas-Fort Worth is about \$5,300 against roughly \$17,500. The
margin on a difference between two metro medians is
\(\mathrm{MOE}_{A-B} = \sqrt{\mathrm{MOE}_A^{2} + \mathrm{MOE}_B^{2}}\),
where \(\mathrm{MOE}_A\) and \(\mathrm{MOE}_B\) are the two 90 percent
margins computed from the replicate weights. The square-root form assumes
the two metro subsamples are independent, which is the usual assumption for
non-overlapping geographies in this file, and it is narrower than adding the
two margins; a gap smaller than that combined margin is one the data cannot
distinguish from zero. Nothing in these three medians supports an assumption
that the Live Music Capital pays its musicians better than the rest of the
state, and nothing in them establishes that it pays them worse.}

\subsection{Austin musicians hold health coverage at roughly the workforce rate}

Health coverage among Austin musicians is \NPumsAustinMusHicov\
(\NPumsAustinMusHicovMoe) against \NPumsAustinAllworkHicov\
(\NPumsAustinAllworkHicovMoe) for all Austin workers, a difference well
inside the musician margin, so the two rates are indistinguishable at the
sample sizes at hand.\footnote{Both shares use the survey's HICOV variable,
which is one when the respondent reports any health insurance coverage
(private, public or both) at the time of the interview, and zero otherwise.
Universe: employed adults aged 18 to 64 in the Austin-Round Rock-San Marcos
MSA, ACS 2020--2024 five-year Texas PUMS; the musician side comes from 77
unweighted records (OCCP 2752), the all-workers side from 48,643. Private and
public coverage are recorded separately (PRIVCOV and PUBCOV) and overlap for
respondents who hold both, so those two shares do not sum to HICOV.}
Private coverage runs the other way: \NPumsAustinMusPrivcov\
(\NPumsAustinMusPrivcovMoe) of Austin musicians hold private insurance
against \NPumsAustinAllworkPrivcov\ for the wider workforce, so the
any-coverage parity rests on public coverage, and possibly on Austin's
music-specific nonprofit health infrastructure, rather than on
employer-sponsored plans.

The 2022 Greater Austin Music Census reports that 84 percent of its
respondents held health coverage in 2022, up from 79 percent in the 2014
census, and its authors attribute that five-point rise to the Health
Alliance for Austin Musicians.\footnote{Sound Music Cities and City of
Austin Music \& Entertainment Division, \emph{2022 Greater Austin Music
Census}, Data Appendix pp.\ 72--73 and Summary Report p.\ 10. Both
comparison figures are self-selected respondent shares rather than
population estimates, and the two censuses are not comparable samples, so
this is a change in reported coverage among the people who answered rather
than a measured change in the workforce as a whole.} The Health Alliance
for Austin Musicians, HAAM, is a 501(c)(3) launched in 2005 that enrols
low-income working Austin musicians (a 400 percent-of-federal-poverty-line
eligibility ceiling) in a coordinated network of primary care, dental,
vision, mental-health and hearing services at subsidised
fees.\footnote{Health Alliance for Austin Musicians, \emph{Annual Impact
Report}, 2005 through 2025 series at
\url{https://www.myhaam.org/annual-reports}. Retrieved 1 August 2026.
Membership rose from 450 in 2005 to 3{,}136 in 2022 and 3{,}149 in 2023;
cumulative healthcare value delivered rose from an initial-year figure of
\$115{,}000 in 2005 to \$46 million (2015), \$73.7 million (2019), \$123
million (2022) and \$215 million (2025), figures drawn from HAAM's own
annual reports. HAAM publishes a member-income share-below-threshold rather
than a member median, so the report does not treat its threshold statistics
as earnings measurements of the wider Austin musician workforce.} SIMS
Foundation, founded in 1995, delivers mental-health and substance-use care
to Austin musicians and their families on a similar
reduced-fee model.\footnote{SIMS Foundation, at
\url{https://simsfoundation.org/about/}, retrieved 1 August 2026.} The ACS
does not identify HAAM or SIMS enrolees, so the coverage rate here is not
causally attributable to either organisation, but the parity is consistent
with what the census reports.

\subsection{Musician households in Texas pay a higher median rent when they rent}

Among Texas households with a working adult, musician households and other
working households rent at close to the same rate: about 41 percent of
musician households are renters, against 39 percent of working households as
a whole.\footnote{Renter shares are weighted-household ratios computed from
the ACS 2020--2024 five-year Texas PUMS housing extract, using weighted
renter and owner household counts, so the two shares are computed on
comparable denominators; the small difference is well inside the musician
margin. Age composition is also similar: the median age of the householder
is close in the two groups, so tenure differences do not reduce to a
younger-therefore-not-yet-owning composition effect at the level the
household file can measure. The Austin metro sub-cell is too small to
publish tenure shares for.} \emph{Among renter households}, musician
households pay a higher median gross rent than working households as a
group.\footnote{The rent statistics come from the ACS housing extract for
Texas (2020--2024 five-year PUMS), a separate file with its own household
weight WGTP and its own 80 replicate weights (WGTP1--WGTP80), used in the
same successive-difference formula as above. Rent is put on 2025 dollars
by ADJHSG (which restates the five survey years on a 2024 basis) and the
same CPI-U 2024-to-2025 factor. Universe: Texas households with at least
one employed resident aged 18 to 64. A household counts as a musician
household when at least one of those employed residents reports OCCP 2752
(Musicians and Singers) as their current or most recent occupation; the
comparison group is every Texas household with at least one employed
resident, so the two groups nest rather than partition, and the comparison
is musician households against all working households, not against a
mutually exclusive remainder. The two sample sizes are 160 unweighted
musician-renter households (representing 4,552 weighted households) and
101,799 unweighted worker-renter households (representing 3.14 million).
Both figures are Texas households rather than Austin ones, because the
Austin-only musician cell is too small in the household file to publish.}
The median monthly gross rent is \NPumsHhMusMedrent\ for musician renter
households against \NPumsHhWorkerMedrent\ for renting Texas worker
households. The rent level is higher, but the burden is not:
\NPumsHhMusRentThreeZero\ (\NPumsHhMusRentThreeZeroMoe) of musician renter
households pay more than 30 percent of household income in
rent\footnote{The rent-burden variable is GRPIP, gross rent as a percentage
of household income, computed by the Bureau from monthly gross rent (GRNTP,
which is contract rent plus tenant-paid utilities) and household income
(HINCP). GRPIP is a ratio of two same-year quantities, so no inflation
adjustment applies; it is top-coded at 101, a category the Bureau also uses
for households reporting zero or negative income. The 30 percent line is the
standard federal cost-burden threshold used by HUD.} against
\NPumsHhWorkerRentThreeZero\ (\NPumsHhWorkerRentThreeZeroMoe) of working
renter households as a whole, a gap well inside the musician margin. Read
together: a musician-renter household in Texas signs a lease on a more
expensive unit than the average working-renter household does, but pays
roughly the same share of household income for it, likely because a second
earner or non-earnings income (a topic Section~\ref{sec:costs} takes up)
makes the rent-to-income ratio hold even as the rent bill runs higher. What
the household file cannot say is why the rent bill runs higher, since the ACS
does not record which parts of the metro a household lives in or which
features of the unit drive its rent.
